% ============================================================================
%  8. Pruebas y resultados
% ============================================================================
\section{Pruebas y resultados}
\label{sec:pruebas}

Las pruebas se hicieron por etapas sobre la tarjeta, de lo más simple a lo más
complejo, para aislar cada falla. La tabla~\ref{tab:resultados} resume las ya
realizadas; las mediciones eléctricas y la demostración con alimentación
definitiva quedan para completar.
% Para llenar: reemplazar cada \pendiente{...} por el valor medido y
% "Pendiente" por "Cumple" o "No cumple".

\subsection{Pruebas funcionales realizadas en la tarjeta}

\begin{table}[H]
  \centering
  \caption{Resultados de las pruebas por etapas.}
  \label{tab:resultados}
  \small
  \begin{tabularx}{\textwidth}{@{}>{\raggedright\arraybackslash}p{3.6cm} L l@{}}
    \toprule
    \textbf{Prueba} & \textbf{Observación} & \textbf{Resultado} \\
    \midrule
    Grabación (tarjeta 1) & El PIC se identifica pero la escritura falla en 0x0000, tras la conexión invertida del ICSP. & No cumple \\
    Grabación (tarjeta 2) & \emph{Programming Successful} con el mismo PICkit y el mismo \texttt{.hex}. & Cumple \\
    Prueba de LEDs (\texttt{prueba\_leds.c}) & Secuencia LED1 $\rightarrow$ LED2 $\rightarrow$ LED3 correcta: PIC, cristal y grabación verificados. & Cumple \\
    Código del docente con motores & Sin movimiento: EEP en \SI{0.14}{\volt}, driver dormido. & No cumple \\
    Enlace Bluetooth (primera versión) & Solo cambiaba LED2: bytes alterados por velocidad equivocada del HC-05. & No cumple \\
    Enlace Bluetooth (con detección de baudios) & LED1 cambia con cada pulsación: tramas de Dabble decodificadas. & Cumple \\
    Comandos sin motores (alimentación del PICkit) & LED3 se enciende con cada dirección y se apaga al soltar; $\triangle$, $\bigcirc$ y $\times$ solo cambian LED1; sin reinicios. & Cumple \\
    Motor con pila de \SI{9}{\volt} & Al arrancar, el HC-05 se desconecta o el PIC se reinicia (tres parpadeos); antes del freno de seguridad, el motor quedaba girando. & No cumple \\
    Motor con EEP en \SI{1.4}{\volt} & Respuesta intermitente; EEP alterna entre \SI{1.4}{\volt} y \SI{5}{\volt}. & No cumple \\
    Motor con EEP a \SI{5}{\volt} y batería 18650 & \pendiente{Por realizar} & Pendiente \\
    \bottomrule
  \end{tabularx}
\end{table}

\subsection{Mediciones de PWM (cristal de 20 MHz)}

\begin{table}[H]
  \centering
  \caption{Mediciones eléctricas esperadas y obtenidas.}
  \label{tab:mediciones}
  \small
  \begin{tabularx}{\textwidth}{@{}L >{\raggedright\arraybackslash}p{3.0cm} c c c@{}}
    \toprule
    \textbf{Prueba} & \textbf{Esperado} & \textbf{Proteus} & \textbf{Hardware} & \textbf{Resultado} \\
    \midrule
    Frecuencia PWM en RC2               & \SI{1250}{\hertz}                 & \pendiente{} & \pendiente{} & Pendiente \\
    Frecuencia PWM en RC1               & \SI{1250}{\hertz}                 & \pendiente{} & \pendiente{} & Pendiente \\
    $T_{on}$ con $DC = 250$             & \SI{200}{\micro\second} (25\,\%) & \pendiente{} & \pendiente{} & Pendiente \\
    $T_{on}$ con $DC = 500$             & \SI{400}{\micro\second} (50\,\%) & \pendiente{} & \pendiente{} & Pendiente \\
    $T_{on}$ con $DC = 750$             & \SI{600}{\micro\second} (75\,\%) & \pendiente{} & \pendiente{} & Pendiente \\
    Duración de rampa 0 $\rightarrow$ 1000 & \SI{2.1}{\second} (21 pasos)  & \pendiente{} & \pendiente{} & Pendiente \\
    Voltaje en VM con motores girando   & $\leq$ \SI{10}{\volt}            & ---          & \pendiente{} & Pendiente \\
    VDD del PIC al arrancar motores     & \SI{5}{\volt} $\pm$ \SI{0.25}{\volt} & ---       & \pendiente{} & Pendiente \\
    EEP (nSLEEP) con motores girando    & \SI{5}{\volt} fijo               & ---          & \pendiente{} & Pendiente \\
    UART: comando \texttt{F} con eco    & Avance y eco \texttt{F}          & \pendiente{} & \pendiente{} & Pendiente \\
    Dabble: cruceta $\uparrow$          & Ambas ruedas en avance, LED3 encendido & ---    & \pendiente{} & Pendiente \\
    \bottomrule
  \end{tabularx}
\end{table}

\subsection{Demostración funcional (rúbrica, criterio 4)}

\begin{table}[H]
  \centering
  \caption{Validación de maniobras en hardware real.}
  \label{tab:demostracion}
  \small
  \begin{tabularx}{\textwidth}{@{}l L L c@{}}
    \toprule
    \textbf{Maniobra} & \textbf{Comportamiento esperado} & \textbf{Observación} & \textbf{Resultado} \\
    \midrule
    Avance         & Ambas ruedas hacia adelante, aceleración suave & \pendiente{} & Pendiente \\
    Reversa        & Ambas ruedas hacia atrás, aceleración suave    & \pendiente{} & Pendiente \\
    Giro izquierda & Rotación antihoraria sobre su eje              & \pendiente{} & Pendiente \\
    Giro derecha   & Rotación horaria sobre su eje                  & \pendiente{} & Pendiente \\
    Freno          & Detención rápida sin reinicio del PIC          & \pendiente{} & Pendiente \\
    \bottomrule
  \end{tabularx}
\end{table}
